\section{Discussion}
\subsection{Implications for Service Orchestration}
The experiments show a practical use for Jev in edge-service admission:
reducing the time and API fees spent interpreting requests while retaining
the observed correct-completion counts. The shared intent contract lets
Jev replace the generative interpreter within the same validation and
scheduling pipeline. The real OCR results show that the latency benefit
remains visible after image transfer, worker queueing, and recognition.

This integration gives service designers two ways to reduce interpretation
overhead. A faster decision backend shortens fresh calls, while caching
reuses decisions for repeated descriptions. The repeated-text conditions
show how effective the latter can be, including for self-hosted Qwen.
Choosing between backends therefore depends on the share of requests that
need a new decision and the time available for service execution. In the
admission path studied here, shorter decision waiting leaves more deadline
slack and releases interpretation capacity sooner.

\subsection{Evaluating the Complete Service Path}
The shared scheduler connects interpreter performance to delivered service.
It reads current worker state after interpretation and checks placement,
tier, priority, and predicted completion time before dispatch. Evaluating
intent fidelity, execution compliance, and actual output separately makes
it possible to trace how an interpretation affects the final result.
For example, the normal-versus-unspecified urgency difference maps to the
same execution priority, while the returned OCR text determines whether
recognition succeeded. This evaluation supports a service-level choice of
interpreter using completion, response time, and fees together.

\subsection{Scope and Next Steps}
The studies use synthetic English requests, repeated texts and images, and
one measurement session; completion retention is descriptive as defined in
Section~\ref{sec:design}. The real service has two workers and one service
family. Hosted timings include provider and network paths, with sequential
runs exposed to temporal variation. They compare deployed services rather
than isolate internal model inference. API fees cover interpretation calls;
local compute and communication costs are outside this measure.
Appendix~\ref{app:settings} records the configuration.

The next step is to evaluate independently authored requests across days,
services, and network conditions, using deadlines drawn from application
requirements. Optimized local serving and a domain-trained classifier or
bounded extractor, absent from the present comparison, would extend the
available deployment choices.
